\fancyhead[L]{\small NOUET'S OBSERVATIONS IN EGYPT.}
\hypertarget{native-539Nouet}{}\label{native-539Nouet}
\begin{center}\small Göttingische gelehrte Anzeigen. No. 79. pp. 780..781. May 18, 1811.\\[3pt]\rule{.72\linewidth}{.3pt}\end{center}

\emph{Description of Egypt. Volume One, printed in Paris by the Imperial Press} 1809. \raisebox{-0.3ex}{,}\emph{First Memoir. Astronomical Observations Made in Egypt During the Years} VI. VII \emph{and} VIII (1798. 1799. 1800) \emph{by} M. \textsc{Nouet} \emph{astronomer of the Commission of Sciences and Arts of Egypt.}

The instruments used by \textsc{Nouet} were a repeating circle 25 centimetres in diameter with a decimal scale (later on a quadrant with a radius of 35 centimetres is also mentioned), a marine timekeeper by \textsc{Louis Berthoud} and an achromatic telescope by \textsc{Dollond}, with an aperture of 63 millimetres. In Cairo he received a second chronometer from \textsc{Beauchamp}, whom he met there. The observations established the latitudes and longitudes of 36 different points. The space available in our journal and its purpose limit how far we can go into the details, so we content ourselves with selecting only a few particulars. The southernmost points are the island of \emph{Philä}, above the cataracts of the Nile, at latitude $24^\circ\,1'\,34''$, and \emph{Syene}, at latitude $24^\circ\,5'\,23''$. There is a well-known story about the latter place, according to which a well was lit to the bottom at the time of the summer solstice. Using \textsc{Laplace}'s formula, \textsc{Nouet} places the time when the obliquity of the ecliptic equalled that latitude at 3430 years before our era. He therefore believes that he must date the story that far back. This conclusion, however, seems insufficiently justified, since such a phenomenon could still have occurred even if the \emph{centre} of the Sun passed a short distance south of the zenith. Most of the longitudes were determined by chronometer, but \textsc{Nouet} also had the opportunity to observe four occultations. He observed the occultation of Venus at Ssalehhiyeh on December 13, 1798, that of $\delta$ in Scorpius at Cairo on April 21, 1799, that of Venus at the same place on November 23, 1799, and that of $\alpha$ in Scorpius at Alexandria on July 28, 1800. He also observed a considerable number of eclipses of Jupiter's satellites at Cairo. The declination of the magnetic needle at Alexandria, as the mean of 26 observations whose dates are not given. was found to be $13^\circ\,6'$ N.\,W., and the inclination, as the mean of 12 observations, $47^\circ\,30'$ north.
